\documentclass[11pt]{article}
\usepackage[T1]{fontenc}
\usepackage{lmodern}
\usepackage[margin=1in]{geometry}
\usepackage{amsmath,amssymb,amsthm,mathtools}
\usepackage{booktabs,microtype}
\usepackage[hidelinks]{hyperref}
\numberwithin{equation}{section}
\newtheorem{theorem}{Theorem}[section]
\newtheorem{proposition}[theorem]{Proposition}
\newtheorem{lemma}[theorem]{Lemma}
\newtheorem{corollary}[theorem]{Corollary}
\theoremstyle{definition}
\newtheorem{definition}[theorem]{Definition}
\theoremstyle{remark}
\newtheorem{remark}[theorem]{Remark}
\newcommand{\HH}{\mathcal H}
\newcommand{\ceil}[1]{\left\lceil #1\right\rceil}
\newcommand{\floor}[1]{\left\lfloor #1\right\rfloor}
\newcommand{\pos}[1]{\left(#1\right)_+}
\newcommand{\bt}{\frac67}
\allowdisplaybreaks[2]
\setlength{\emergencystretch}{2em}
\title{A six-sevenths bound for transversals in hypergraphs with property $(7,2)$}
\author{}
\date{}
\begin{document}
\maketitle
\begin{abstract}
A hypergraph has property $(7,2)$ if every subfamily of at most seven
edges has a transversal of size at most two. We prove that every
$k$-uniform hypergraph with this property satisfies
$\tau(\HH)\le\lceil6k/7\rceil+1$ for $k\ge7$.
This improves the previously established asymptotic upper coefficient
$7/8$ to $6/7$. The proof combines sharper local avoidance constructions
with three successive exclusions of pair-intersection sizes and a
maximal-intersection argument. All constructions admit integral
allocations of points among at most four requests.
\end{abstract}

\section{Introduction}
A \emph{transversal} of a hypergraph is a set meeting every edge, and its
\emph{transversal number} $\tau(\HH)$ is the minimum size of such a set.
For integers $k\ge2$ and $r\ge3$, let $f(k,r)$ be the largest transversal
number of a finite $k$-uniform hypergraph in which every subfamily of at
most $r$ edges has a transversal of size at most two. The use of ``at most''
avoids vacuous exceptions for families with fewer than $r$ edges.

Erd\H{o}s, Fon-Der-Flaass, Kostochka and Tuza studied these local
transversal conditions in~\cite{EFKT}. For $r=7$, Fon-Der-Flaass,
Kostochka and Woodall proved the bounds
\[
 \ceil{3k/4}\le f(k,7)\le\ceil{7k/8}
\]
for the upper-bound range $k\ge8$; see~\cite{FKW}.
The question whether $f(k,7)=(3/4+o(1))k$ remains the motivating problem.
Here we improve the coefficient in the upper bound to $6/7$.

\begin{theorem}\label{thm:main}
Let $k\ge7$. Every finite $k$-uniform hypergraph with property $(7,2)$
satisfies
\[
 \tau(\HH)\le\ceil{6k/7}+1.
\]
Consequently $\limsup_{k\to\infty}f(k,7)/k\le6/7$.
\end{theorem}

Our proof refines the avoidance method of Fon-Der-Flaass, Kostochka
and Woodall~\cite{FKW}, which already uses good triples, intersection
gaps, and a maximal small intersection. The improvement comes from
stronger local completion constructions and their three-stage
combination. Broader local-cover problems were subsequently studied
by Kostochka~\cite{Kostochka} and Buci\'c, Kor\'andi and Sudakov~\cite{BKS}.

The proof uses the elementary consequence of $\tau(\HH)>T$ that each
set of at most $T$ points is avoided by an edge of $\HH$. Starting from
three edges with empty common intersection, we construct four such
requests so that no pair of points can meet all seven resulting edges.
Some constructions use a previously established gap in the possible
pair-intersection sizes. This permits three successive exclusions,
after which a pair with largest intersection at most $k/2$ supplies the
final contradiction.

The main argument is given in Sections~\ref{sec:gaps} and~\ref{sec:finish}.
The elementary setup is in Section~\ref{sec:prelim}. For completeness,
Appendices~\ref{app:allocation}--\ref{app:smallmax} contain full proofs of
all local constructions used in the argument. The additive constant is
handled within these constructions; a single budget
$T=\lceil6k/7\rceil+1$ is used throughout.

\section{Avoidance requests and good triples}\label{sec:prelim}
Throughout the main argument, $\HH$ is a finite $k$-uniform hypergraph
with property $(7,2)$ and $\tau(\HH)>T$, where $T$ is an integer. An \emph{avoidance request} is a set $D$ of at most $T$ points.
There is an edge disjoint from $D$, called a response to the request.
We construct responses successively. The requests need not be disjoint,
and their sizes are bounded individually.

A triple $E,F,G$ is \emph{good} if $E\cap F\cap G=\varnothing$.
Write
\[
 X=E\cap F,\qquad Y=E\cap G,\qquad Z=F\cap G,
\]
and write $P_E,P_F,P_G$ for the private parts of its three edges.
The six cells $X,Y,Z,P_E,P_F,P_G$ are disjoint. Any pair meeting the
three edges either lies in two distinct pair cells or lies in a pair
cell and the opposite private part. Points outside $E\cup F\cup G$
are irrelevant, because their partner alone cannot meet a good triple.
This observation verifies the static constructions below.

For four edges with all triple intersections empty, a pair meeting all
four must lie in two complementary pair cells. If precisely one triple
cell survives, there is in addition the product of that cell with the
opposite edge. These descriptions will be used explicitly in the local
proofs.

\begin{lemma}[Balanced request]\label{lem:balanced}
Suppose $E,F$ have intersection size $M\le T\le k$. There is a request
of size $T$ containing $E\cap F$ such that an avoiding response $G$
forms a good triple and
\[
 |E\cap G|,|F\cap G|\le k-\floor{(T+M)/2}.
\]
\end{lemma}
\begin{proof}
Include the $M$ common points and distribute the remaining $T-M$
requested points as evenly as possible between $E\setminus F$ and
$F\setminus E$. Both selections fit because $T\le k$. Each new trace
has size at most $k-M-\lfloor(T-M)/2\rfloor$, which is the asserted
quantity.
\end{proof}

Most scalar inequalities will be written after division by $k$.
Thus lower-case pair sizes in Sections~\ref{sec:gaps} and~\ref{sec:finish}
are proportions. The appendices use actual integer sizes unless stated
otherwise. A triple \emph{closes at budget $T$} if four further requests
of size at most $T$ force a subfamily of at most seven edges with no
transversal of size at most two. Each closing result is therefore
incompatible with property $(7,2)$ under the assumption $\tau(\HH)>T$.

Throughout the main proof put
\begin{equation}\label{eq:budget}
 \beta=\frac67,\qquad T=\ceil{\beta k}+1.
\end{equation}
For $k\ge7$ we have $T\le k$. Every local result below is implemented
with this same integer budget. In particular, the losses in different
requests are not added together.

\section{Three gaps in the intersection spectrum}\label{sec:gaps}
We first explain the initial request used to obtain each gap. Suppose
$|E\cap F|=qk$. Choose $u$ and put $v(q)=2-\beta-q-u$.
If $0\le u,v(q)\le1-q$, retain subsets of the private parts of $E,F$
of sizes $\lfloor uk\rfloor$ and $\lfloor v(q)k\rfloor$, and request
all remaining points of $E\cup F$. The request contains $E\cap F$ and
has size
\[
 2k-qk-\floor{uk}-\floor{v(q)k}<\beta k+2.
\]
Its size is an integer and is therefore at most $T$. The response gives
a good triple whose other two pair proportions are at most $u,v(q)$.

The choices needed are displayed in Table~\ref{tab:gaps}. In each row,
the host conditions $0\le u,v(q)\le1-q$ hold throughout the stated
interval. In the second and third rows, gaps already proved reduce the
raw response bounds to the bounds in the final column.
\begin{table}[ht]
\centering
\begin{tabular}{ccccc}
\toprule
Stage & Initial $q$ & $u$ & $v(q)$ & Response bounds\\
\midrule
1 & $[3/7,10/21]$ & $5/14$ & $11/14-q$ & $y,z\le5/14$\\
2 & $[4/21,5/14]$ & $10/21$ & $2/3-q$ & $y,z<3/7$\\
3 & $[10/21,1/2]$ & $1/3$ & $17/21-q$ & $y,z<4/21$\\
\bottomrule
\end{tabular}
\caption{Successive exclusions of pair-intersection sizes.}\label{tab:gaps}
\end{table}

\subsection{The first gap}
\begin{proposition}\label{prop:gapone}
Under \eqref{eq:budget} and $\tau(\HH)>T$, no pair of edges has
intersection proportion in $[3/7,10/21]$.
\end{proposition}
\begin{proof}
The first row of Table~\ref{tab:gaps} reduces the assertion to closing
a good triple with
\[
 3/7\le x\le10/21,\qquad 0\le z\le y\le5/14.
\]
Put $s=y+z$, $d=y-z$, and $S=x+s$.

\paragraph{Case 1: $d>x-1/7$.}
Use Lemma~\ref{lem:asymone} in orientation $(x,y,z)$.
Since $s+d=2y\le5/7$, its four bounds are verified by
\begin{align*}
 S&<6/7,& 1/2+y&\le6/7,\\
 (1+2x-d)/2&<(8/7+x)/2\le17/21,\\
 (1+2x+2s-d)/3&<(20/7-x)/3\le17/21.
\end{align*}

\paragraph{Case 2: $d\le x-1/7$ and $s\le11/7-2x$.}
Use Lemma~\ref{lem:fourcase} in orientation $(y,z,x)$.
Its nine terms satisfy
\begin{align*}
 s&\le5/7,&1/2+y, 1/2+z&\le6/7,\\
 1+d-x&\le6/7,&1-d-x&\le4/7,\\
 1-S/3&\le6/7,&(3+S)/5&\le88/105,\\
 (1+s+2x)/3&\le6/7,&(2+3x)/4&\le6/7.
\end{align*}
Here $S\ge x\ge3/7$ and $S\le25/21$.

\paragraph{Case 3: $s>11/7-2x$.}
Use the symmetric allocation of Lemma~\ref{lem:symmetric}.
Since $x\le10/21$,
\[
 s>x+1/7,\qquad x+y-z=x+d<3x-6/7\le4/7,\qquad S\le25/21<9/7.
\]
The second inequality uses $d\le5/7-s$. These are precisely the three
hypotheses of that lemma. Its integral allocation fits $T$ with no
rounding loss. The three cases exhaust the domain.
\end{proof}

\subsection{The middle gap}
\begin{proposition}\label{prop:gaptwo}
Under the same assumptions, no pair has intersection proportion in
$[4/21,5/14]$.
\end{proposition}
\begin{proof}
The second row of Table~\ref{tab:gaps} gives two new traces of proportion
at most $10/21$. Proposition~\ref{prop:gapone} forces both below $3/7$.
It suffices to close a good triple with
\[
 4/21\le x\le5/14,\qquad 0\le z\le y\le3/7.
\]
If $x+z\le11/21$, then $z\le1/3<5/14$. Apply
Lemma~\ref{lem:hallstatic} with distinguished pair size $a=y$ and
remaining sizes $b=x,c=z$: its conditions follow from
$a\le3/7$, $b,c\le5/14$ and $b+c\le11/21<4/7$.

If $x+z\ge11/21$, use Lemma~\ref{lem:gaptwo} with orientation
$(y,x,z)$ and the established gap $[\ell,h]=[3/7,10/21]$.
Its domain conditions hold because $y\ge z$ and
$x+z\ge11/21=1-h$. Its six budget terms are bounded by
\begin{align*}
 2-2h-x&=22/21-x\le6/7,\\
 y+\ell, z+\ell&\le6/7, &1/2+x&\le6/7,\\
 3/4&<6/7, &(3+x+y+z)/5&\le59/70<6/7.
\end{align*}
The lemma's integer construction uses $T=\lceil\beta k\rceil+1$.
\end{proof}

\subsection{The final gap}
\begin{proposition}\label{prop:gapthree}
Under the same assumptions, no pair has intersection proportion in
$[10/21,1/2]$.
\end{proposition}
\begin{proof}
The last row of Table~\ref{tab:gaps} gives new traces at most $1/3$.
Proposition~\ref{prop:gaptwo} forces both below $4/21$. We therefore
consider a good triple with
\[
 10/21\le x\le1/2,\qquad 0\le z\le y\le4/21.
\]
Write $S=x+y+z$, and set $\ell=4/21$, $h=5/14$.

\paragraph{Case 1: $S\le5/7$.}
Lemma~\ref{lem:split} applies, since
\[
 (1+S)/2\le6/7,\qquad 1-x+|y-z|\le5/7,\qquad 1/3+x\le5/6.
\]

\paragraph{Case 2: $5/7\le S\le6/7$ and $z\le1/14$.}
Apply Lemma~\ref{lem:asymtwo} in orientation $(x,z,y)$.
Besides $S\le\beta$, its bounds follow from
\[
 1/3+x\le5/6,\qquad 1-x+y\le5/7,\qquad
 2x+y+3z\le59/42<11/7.
\]

\paragraph{Case 3: $5/7\le S\le6/7$ and $z\ge1/14$.}
Apply Lemma~\ref{lem:gapempty} with the gap $[\ell,h]$.
The nine required inequalities follow from
\begin{align*}
 S&\le6/7,&1-h+y, 1-h+z&\le5/6,\\
 2-2h-x&\le17/21,&1/3+x&\le5/6,\\
 2x-y, 2x-z&\le13/14,&x-y-z&=2x-S\le2/7,\\
 2\ell+y+z&\le16/21.
\end{align*}
In particular, the two terms with numerator $(2-h)+2x-y$ or
$(2-h)+2x-z$ are at most $(23/14+13/14)/3=6/7$, and the term with
numerator $3-2h+x-y-z$ is at most $(16/7+2/7)/3=6/7$.

\paragraph{Case 4: $S\ge6/7$.}
Apply Lemma~\ref{lem:gapsurvive}. Its domain condition is this case,
and its six bounds follow from
\begin{align*}
 x+y&\le29/42,&1-h+y, 1-h+z&\le5/6,\\
 1/4+x+(y+z)/4&\le71/84,&2\ell+y+z&\le16/21,\\
 1/2+x/2+(z+\ell)/4&\le71/84.
\end{align*}
This exhausts the domain.
\end{proof}

Propositions~\ref{prop:gapone} and~\ref{prop:gapthree} together show that
there is no pair intersection in $[3k/7,k/2]$. We now prove that the
weaker condition ``at most $3k/7$ or greater than $k/2$'' already suffices
for a slightly stronger bound, with no additive allowance when $k\ge8$.

\section{The maximal-intersection argument}\label{sec:finish}
\begin{proposition}\label{prop:localfinish}
Suppose every pair intersection in $\HH$ is at most $mk$ or greater
than $k/2$. A good triple with sorted pair proportions $m\ge y\ge z$
whose largest pair intersection is $mk$ closes at integer budget
$T_0=\lceil6k/7\rceil$ whenever
\begin{equation}\label{eq:finishdomain}
 0\le m\le3/7.
\end{equation}
\end{proposition}
\begin{proof}
In this proof put $T=T_0$ and $S=m+y+z$. In particular
$S\le3m\le9/7$.

\paragraph{Case 1: $y\le5/14$.}
If $y-z>m-1/7$, apply Lemma~\ref{lem:asymone} in orientation
$(m,y,z)$. Its terms satisfy
\begin{align*}
 S&<2y+1/7\le6/7,&1/2+y&\le6/7,\\
 (1+2m-y+z)/2&<(1+m+1/7)/2\le11/14,\\
 (1+2m+y+3z)/3&<(10/7+4y-m)/3\le5/6.
\end{align*}
If $y-z\le m-1/7$ and $S<3/7$, use the same lemma in orientation
$(z,y,m)$. Its last two terms are at most $(1+S)/2<5/7$ and
$(1+3S)/3<16/21$, respectively.

In the remaining subcase, $y-z\le m-1/7$ implies $m\ge1/7$.
Use Lemma~\ref{lem:fourcase} in orientation $(z,y,m)$. Its nine bounds
follow from
\begin{gather*}
 y+z\le5/7,\quad 1/2+y, 1/2+z\le6/7,\quad
 m+y-z\ge m\ge1/7,\quad m-y+z\ge1/7,\\
 S\ge3/7,\qquad (3+S)/5\le29/35,\qquad
 (1+y+z+2m)/3\le6/7,\qquad (2+3m)/4\le23/28.
\end{gather*}

\paragraph{Case 2: $y>5/14$ and $1/7<z\le5/14$.}
Use the full-gap construction of Lemma~\ref{lem:gaptwo}, with orientation
$(m,z,y)$, $h=1/2$ and $\ell=m$. This version is valid when intersections
are at most $mk$ or greater than $k/2$: its proof uses only the implication
that an intersection at most $k/2$ is at most $mk$. The two domain
conditions follow from $m+z\ge y+z>1/2$. We check its initial integer
request separately to remove the additive allowance in that lemma.
Write $Z=zk$ for the integer pair size used in its middle coordinate.
The exact initial request size is
\[
 Z+2(k-Z-\lfloor k/2\rfloor)_+
   =\max\{Z,k-Z+(k\bmod2)\}.
\]
Since $Z>k/7$, integrality gives $Z\ge\lfloor k/7\rfloor+1$, so
this request fits $T=k-\lfloor k/7\rfloor$. The domain inequalities
also ensure that the cuts fit their hosts when $k$ is odd. All subsequent
allocations in the cited proof are integral. Its six homogeneous bounds
are
\[
 \begin{gathered}
 1-z\le6/7,\quad 2m\le6/7,\quad y+m\le6/7,\\
 1/2+z\le6/7,\quad 3/4<6/7,\quad (3+S)/5\le59/70.
 \end{gathered}
\]
In this version, the relevant strict trace bounds in the lemma's proof
become weak bounds; all capacity estimates remain valid.

\paragraph{Case 3: $y>5/14$ and $z\ge5/14$.}
Apply Lemma~\ref{lem:symmetric}. Indeed,
\[
 y+z>5/7>m+1/7,\qquad S\le3m\le9/7.
\]
Moreover $m+y-z\le2m-z\le6/7-5/14=1/2<4/7$.

\paragraph{Case 4: $y>5/14$ and $z\le1/7$.}
Use Lemma~\ref{lem:nearcore} in orientation $(x,y,z)=(m,z,y)$.
The full-gap parameter in that lemma is still $m$, and
$5/14<y\le m\le3/7$.
The domain condition and the first five bounds of the lemma follow from
\begin{align*}
 z+y&\le4/7,&m+y&\le6/7,\\
 1+z-y&<11/14,&2+m-2y&\le17/7-2y<12/7,\\
 2+2m-y&<5/2<18/7,&2m&\le6/7.
\end{align*}
If $S\le6/7$, use its empty-core variant. The two remaining bounds are
\[
 1-m+z<11/14,\qquad 2-2m+y\le2-m<23/14<12/7.
\]
If $S\ge6/7$, put $\Delta=S-6/7$. The four remaining bounds, after
substituting this value, become
\begin{gather*}
 1+2z+y\le12/7,\qquad 2-m+z+2y\le18/7,\\
 2+m+z\le18/7,\qquad 3+y\le24/7.
\end{gather*}
They follow from
\[
 1+2z+y\le12/7,\qquad 2-m+z+2y\le2+y+z\le18/7,
\]
$m+z\le4/7$, and $3+y\le24/7$.
All these inequalities were checked at the real budget $6k/7$.
Increasing it to the integer $T$ decreases the actual $\Delta$ and
increases every capacity. The integral construction of
Lemma~\ref{lem:nearcore} therefore applies.
\end{proof}

\begin{theorem}[Conditional finishing theorem]\label{thm:finish}
Let $k\ge8$. If a $k$-uniform hypergraph has property $(7,2)$ and every
pair intersection is at most $3k/7$ or greater than $k/2$, then
$\tau(\HH)\le\lceil6k/7\rceil$. The same hypotheses give
$\tau(\HH)\le7$ when $k=7$.
\end{theorem}
\begin{proof}
For $k\ge8$ set $T=\lceil6k/7\rceil=k-\lfloor k/7\rfloor$;
for $k=7$ set $T=7$. Suppose $\tau(\HH)>T$.
An edge avoiding a $T$-subset of another edge gives a pair intersection
at most $k-T<k/2$. Therefore a largest pair intersection $M\le k/2$
exists, and the hypothesis implies $M\le3k/7$.

If $M\le2k/7$, use Lemma~\ref{lem:smallmax}. Its four arithmetic
conditions hold for this $T$. For $k\ge14$, the bounds $T\ge6k/7$
give respective margins at least $k/7-1$, $k/14-1/2$,
$3k/14-1/2$, and $2k/7-4$, with the required strictness in the first.
For $8\le k\le13$, $T=k-1$, and the four conditions reduce to
$k>7$, $k-3\ge\lceil k/2\rceil$, $k\ge5$, and $k\ge8$.
They also hold directly for $k=T=7$.
Otherwise choose a pair attaining $M$ and apply
Lemma~\ref{lem:balanced}. The other two intersections $Y\ge Z$ satisfy
\[
 Y,Z\le k-\floor{(T+M)/2}.
\]
Also $T\ge6k/7$ and $M>2k/7$ give
\[
 Y\le k-\floor{(T+M)/2}<3k/7+1/2\le k/2
\]
for $k\ge7$. By maximality, $Y,Z\le M$. Dividing by $k$ gives the hypotheses of
Proposition~\ref{prop:localfinish}, which produces the required
contradiction. For $k=7$, the local proposition closes at budget
$T_0=6$, and its requests are also available under $\tau(\HH)>7$.
\end{proof}

\begin{proof}[Proof of Theorem~\ref{thm:main}]
Suppose $\tau(\HH)>T$. The three propositions in
Section~\ref{sec:gaps} exclude all pair intersections in $[3k/7,k/2]$.
Theorem~\ref{thm:finish} then contradicts $\tau(\HH)>T$.
\end{proof}

\begin{corollary}
The conclusion of Theorem~\ref{thm:main} also holds for finite families
of nonempty sets of rank at most $k$.
\end{corollary}
\begin{proof}
Pad each edge to size $k$ with points private to that edge. Old
transversals still meet the padded edges. Conversely, replace each
private point of a transversal by an arbitrary point of its original
nonempty edge. This gives a transversal of the original family of no
larger size. Thus padding preserves the transversal number, and it
preserves property $(7,2)$ because every original two-point transversal
still works. Apply Theorem~\ref{thm:main} to the padded family.
\end{proof}

\begin{corollary}\label{cor:arithmetic}
For every integer $n\ge1$,
\[
 f(42n,7)\le36n,\qquad f(42n-1,7)\le36n.
\]
Thus the additive point in Theorem~\ref{thm:main} can be omitted when
$k\equiv0$ or $41\pmod{42}$.
\end{corollary}
\begin{proof}
Set $k=42n$ and $T=6k/7=36n$. Use the same three gap stages as in
Section~\ref{sec:gaps}. All fixed caps and gap thresholds used there
are integral multiples of a point:
\[
 5k/14=15n,\quad 10k/21=20n,\quad k/3=14n,\quad
 4k/21=8n,\quad 3k/7=18n,\quad k/2=21n.
\]
The second cap in each initial request is also integral, since it is
$2k-T-qk$ minus the first cap. Hence those requests have size exactly
$T$, with no floor loss.

Lemmas~\ref{lem:split}, \ref{lem:fourcase}, \ref{lem:asymone},
\ref{lem:asymtwo}, \ref{lem:hallstatic}, and \ref{lem:symmetric} already
have integral constructions at every integer budget satisfying their
homogeneous bounds. The three gap-dependent constructions also fit this
budget. In Lemma~\ref{lem:gapempty}, $hk$ is integral, so the initial
request and the numerator of its split bound incur no rounding error.
That numerator is at most $3T$, and each split request therefore has
size at most $\lceil(3T-T)/2\rceil=T$. The third request is smaller
than $T$. In Lemma~\ref{lem:gapsurvive}, the first two bases are integer
sets of size at most $T+1/2$, hence at most $T$; the total-load argument
then uses integral residual capacities. In Lemma~\ref{lem:gaptwo},
$hk$ is again integral, so the first request fits $T$ exactly, and all
later allocations are integral.

Thus all three gap exclusions hold with this smaller budget.
Theorem~\ref{thm:finish} completes the proof for $k=42n$.
For $k=42n-1$, privately pad to rank $42n$ and apply the first assertion.
The resulting bound $36n$ equals $\lceil6(42n-1)/7\rceil$.
\end{proof}

\appendix
\section{Integral allocations and a near-core construction}\label{app:allocation}
We use the following elementary form of integral capacitated Hall's
theorem. Suppose disjoint sets of points may each be assigned to a
specified set of request indices, and each request has an integer
remaining capacity. An assignment exists if and only if every collection
of source sets has total size at most the sum of the capacities of its
available request indices. Indeed, make a bipartite flow network, with
source-set sizes and request capacities on its outer arcs. The min-cut
conditions are exactly these inequalities; integral augmenting paths
give an assignment of individual points. We will list every nonredundant
condition when this fact is used.

\begin{lemma}[Exact boxed-triangle allocation]\label{lem:triangle}
Let $M_1,M_2,M_3$ be nonnegative integers and let $L_1,L_2,L_3$ be
integers satisfying $L_i\le M_j+M_l$ whenever $\{i,j,l\}=\{1,2,3\}$.
The minimum of $a_1+a_2+a_3$ over integral vectors satisfying
\[
 0\le a_i\le M_i,\qquad a_j+a_l\ge L_i
\]
is the ceiling of
\begin{equation}\label{eq:triangleexact}
 \max\left\{0,L_1,L_2,L_3,
 L_1+L_2-M_3,L_1+L_3-M_2,L_2+L_3-M_1,
 \frac{L_1+L_2+L_3}{2}\right\}.
\end{equation}
\end{lemma}
\begin{proof}
For a proposed integer total $t$, the conditions are equivalent to
\[
 0\le a_i\le\min(M_i,t-L_i),\qquad \sum_i a_i=t.
\]
Integral coordinates exist if and only if all three upper bounds are
nonnegative and their sum is at least $t$. Expand the sum of the three
minima as the minimum of its eight affine expressions. The resulting
conditions are
\begin{gather*}
 t\ge0,L_1,L_2,L_3,\qquad
 t\ge L_i+L_j-M_l\quad(\{i,j,l\}=\{1,2,3\}),\\
 2t\ge L_1+L_2+L_3,\qquad t\le M_1+M_2+M_3,\qquad
 L_i\le M_j+M_l.
\end{gather*}
The last three host inequalities hold by assumption. Every lower bound
in \eqref{eq:triangleexact} is at most $M_1+M_2+M_3$, so the ceiling of
the maximum also satisfies the upper bound on $t$. Coordinates of this
total can be filled successively within their integral upper bounds.
The same argument without ceilings gives the real minimum.
\end{proof}

\begin{lemma}[Exact symmetric static allocation]\label{lem:symmetric}
Let a good triple have pair sizes $m\ge y\ge z\ge0$, and put $S=m+y+z$.
The four-request construction below can be implemented at integer
budget $T$ if and only if
\begin{equation}\label{eq:symexact}
 T\ge\max\left\{k+m-y-z,\frac{k+m}{2},
             \frac{2k+m+y-z}{3},\frac{3k+S}{5}\right\}.
\end{equation}
Whenever it can be implemented, it closes the triple.
In particular, at unit scale, the sufficient conditions
\begin{equation}\label{eq:symconvenient}
 y+z\ge m+1-\beta,\qquad m+y-z\le3\beta-2,\qquad
 m+y+z\le5\beta-3
\end{equation}
with $0\le\beta\le1$ imply closure at every integer budget
$T\ge\beta k$.
\end{lemma}
\begin{proof}
Split the pair cells $M,Y,Z$ into selected parts of sizes $a,b,c$ and
their complements. With the private cells denoted by $P_E,P_F,P_G$,
use four requests
\begin{align*}
 R_0&=M_1\cup Y_1\cup Z_1,\\
 R_M&=M\cup P_G\cup Y_2\cup Z_2,\\
 R_Y&=Y\cup P_F\cup M_2\cup Z_2,\\
 R_Z&=Z\cup P_E\cup M_2\cup Y_2.
\end{align*}
Their sizes are $a+b+c$, $k+m-b-c$, $k+y-a-c$, and $k+z-a-b$.
Apply Lemma~\ref{lem:triangle} with hosts $(m,y,z)$ and requirements
\[
 L_1=k+m-T,\qquad L_2=k+y-T,\qquad L_3=k+z-T.
\]
The host conditions reduce to $T\ge k+m-y-z$. The individual
requirements $L_i\le T$ reduce to $2T\ge k+m$; the two-$L$ requirements
reduce to $3T\ge2k+m+y-z$; and the total requirement is
$5T\ge3k+S$. This proves the exact criterion.

Every pair meeting the good triple is contained in a request. A pair
cell and its opposite private part lie in the corresponding request.
For two different pair cells, if both points lie in selected parts they
belong to $R_0$; otherwise one of the last three requests contains both.
Thus four responses close the triple.

Finally, \eqref{eq:symconvenient} directly bounds the first, third, and
fourth forms in \eqref{eq:symexact} by $\beta$ at unit scale. Its first
and third inequalities imply $2m+1-\beta\le5\beta-3$, so
$m\le3\beta-2\le2\beta-1$. The second form is therefore at most
$\beta$ as well. The exact criterion supplies integral splits without
any rounding allowance. Its optimality concerns this specified
construction only, not all possible closing constructions.
\end{proof}

\begin{lemma}[A static Hall allocation]\label{lem:hallstatic}
A good triple with pair proportions $a,b,c$ closes at every integer
budget $T\ge6k/7$ if
\[
 a\le3/7,\qquad b,c\le5/14,\qquad b+c\le4/7.
\]
\end{lemma}
\begin{proof}
We describe the allocation at unit scale with budget $t=6/7$.
Let the pair cells be $X=E\cap F$, $Y=E\cap G$, $Z=F\cap G$,
of sizes $a,b,c$. Give them request-index labels
$\{1,2,3\}$, $\{2,4\}$, and $\{3,4\}$, respectively.
A label lists all requests that contain its point. Assign points of the
private cells as follows:
\[
 P_E:\{3\}\text{ or }\{4\},\qquad
 P_F:\{2\}\text{ or }\{4\},\qquad
 P_G:\{1\},\{2\}\text{ or }\{3\}.
\]
The four remaining capacities are
\[
 t-a,\quad t-a-b,\quad t-a-c,\quad t-b-c,
\]
which are nonnegative. The seven Hall inequalities for the three
private cells are
\begin{gather*}
 1+2c\le2t,\qquad 1+2b\le2t,\qquad 1+3a\le3t,\\
 2+b+c\le3t,\qquad 2+2a+c\le4t,\qquad
 2+2a+b\le4t,\qquad 3+a\le4t.
\end{gather*}
The first four and the last follow immediately from the stated bounds;
the fifth and sixth have left side at most
$2+6/7+5/14=45/14<24/7$.
At integer scale, the same inequalities hold for every integer
$T\ge6k/7$, and integral flow allocates the private points.

Every pair meeting the good triple has intersecting labels. The three
pair-cell labels intersect pairwise. The allowed labels of each private
cell meet the label of its opposite pair cell. Hence each candidate
pair is contained in one request and is missed by its response.
\end{proof}

\begin{lemma}[Near-core allocation]\label{lem:nearcore}
Suppose every pair intersection in the family is at most $m$ or greater
than $k/2$, where $m\le k/2$. Let $E,F,G$ be a good triple with pair
sizes $x=|E\cap F|$, $y=|E\cap G|$, $z=|F\cap G|$, and put
$\Delta=(x+y+z-T)_+$. The triple closes at integer budget $T$ if
$y+z\le T$ and
\begin{align}
 m+z&\le T,& k+y-z&\le T,\label{eq:ncfirst}\\
 2k+x-2z&\le2T,&2k+x+m-z&\le3T,\nonumber\\
 x+m&\le T,&&\nonumber\\
 k-x+y+\Delta&\le T,&2k-2x+z+\Delta&\le2T,\label{eq:nclast}\\
 2k-z+\Delta&\le2T,&3k-x-y+\Delta&\le3T.\nonumber
\end{align}
If $x+y+z\le T$, the last four bounds may instead be replaced by
\begin{equation}\label{eq:emptycore}
 k-x+y\le T,\qquad 2k-2x+z\le2T.
\end{equation}
\end{lemma}
\begin{proof}
Put $X=E\cap F$, $Y=E\cap G$, $Z=F\cap G$. Request an edge $H$
avoiding $Y\cup Z$ and a subset of $X$ of size
$\min(x,T-y-z)$. The only possibly nonempty triple cell is
$P=E\cap F\cap H$, of size $p\le\Delta$.
Define the disjoint cells
\begin{align*}
 X'&=X\setminus P, & A&=H\cap\bigl(E\setminus(F\cup G)\bigr),\\
 B&=H\cap\bigl(F\setminus(E\cup G)\bigr), &C&=H\cap G,\\
 W&=G\setminus(E\cup F\cup H).
\end{align*}
Write $a,b,c,w$ for the latter four sizes. Then
\begin{gather*}
 a\le k-x-y,\quad b\le k-x-z,\quad c\le k-y-z,\\
 p+a+b+c\le k,\qquad c+w=k-y-z.
\end{gather*}
The complete candidate products for pairs meeting $E,F,G,H$ are
\[
 P\times(Y\cup Z\cup C\cup W),\qquad
 X'\times C,\qquad Y\times B,\qquad Z\times A.
\]
A point in three edges lies in $P$; otherwise a covering pair must lie
in complementary pair cells. This proves completeness of the list.

\paragraph{Case 1: $p+a\le m$.}
Put $P$ in all three remaining requests, $X'$ in requests $1,3$,
$Y,B$ in request $1$, and $Z,A$ in request $2$. Partition $C$ between
requests $1,3$, and then partition $W$ among all three. The initial
loads are
\[
 L_1=x+y+b,\qquad L_2=p+z+a,\qquad L_3=x.
\]
They fit by $L_1\le k+y-z$, $L_2\le m+z$, and $L_3\le x+m$.
The first partition fits because
\[
 L_1+L_3+c=2x+y+b+c\le2k+x-2z\le2T.
\]
After that partition, the total load including $W$ is
\[
 L_1+L_2+L_3+c+w=2x+k+p+a+b\le2k+x+m-z\le3T.
\]
All capacities are integral, so the partitions exist. Each candidate
product is covered: $P$ lies in every request; the two possible labels
of $C$ meet the label of $X'$; and $Y,B$ and $Z,A$ share requests.

\paragraph{Case 2: $p+a>m$.}
The global dichotomy gives $p+a>k/2$. The traces $E\cap H$ and
$G\cap H$ are disjoint, so $c<k/2$ and hence $c\le m$.
Put $P$ in requests $1,2$, $X',C$ in request $2$, $Y,B$ in request
$1$, and $Z$ in requests $1,3$. Partition $A$ between requests $1,3$
and $W$ between requests $1,2$. The initial loads are
\[
 L_1=p+y+z+b,\qquad L_2=x+c,\qquad L_3=z.
\]
They fit by $L_1\le\Delta+k-x+y$, $L_2\le x+m$, and $L_3\le m+z$.
The three Hall conditions for the two flexible cells are
\[
 a\le2T-L_1-L_3,\qquad w\le2T-L_1-L_2,\qquad
 a+w\le3T-L_1-L_2-L_3.
\]
They follow from
\begin{align*}
 L_1+L_3+a&\le\Delta+2k-2x+z\le2T,\\
 L_1+L_2+w&=p+x+k+b\le\Delta+2k-z\le2T,\\
 L_1+L_2+L_3+a+w&=p+x+k+z+a+b
                  \le\Delta+3k-x-y\le3T.
\end{align*}
Thus an integral allocation exists. The labels of $P$ meet all allowed
labels of $Y,Z,C,W$; $X',C$ share request $2$; $Y,B$ share request $1$;
and both labels of $A$ meet the label $\{1,3\}$ of $Z$. Every candidate
pair is therefore missed by a response.

Finally, if $x+y+z\le T$, the first request omits the whole pair-cell
union, so $P=\varnothing$. The cell $W$ has no candidate-pair neighbour
and can be left out of every remaining request. In Case 2 only the
allocation of $A$ between requests $1,3$ is needed. Its conditions are
exactly the base-capacity bounds and \eqref{eq:emptycore}. Case 1 already
works even if $W$ is included. This proves the empty-core variant.
\end{proof}

% Complete local avoidance constructions. All mathematical statements
% are independent of the source-note numbering.
\section{Local avoidance constructions}\label{app:local}
Throughout this appendix, the family is $k$-uniform and has transversal
number greater than the integer request budget $T\le k$. Thus every set
of at most $T$ points is avoided by a family edge. All edges obtained by
such requests belong to the original family. Repetitions cause no problem:
the contradiction is a subfamily of at most seven edges.

The notation of a good triple is as in Section~\ref{sec:prelim}.
A construction \emph{closes} a triple if four further requests force a
subfamily with no transversal of size at most two. Pair-cell sizes in
this appendix are unnormalised integers. Whenever gap parameters and a
real budget coefficient occur, assume $0\le\ell\le h\le1$ and
$0\le\beta\le1$.

\begin{lemma}[Three small intersections; cf.~\cite{FKW}]\label{lem:smalltriple}
Suppose all three pairwise intersections of a good triple have size at most $m\le k/2$. It extends to a bad subfamily of at most seven edges whenever
\[
\tau>B:=\left\lceil\max\{(3k+m)/4,(2k+2m)/3\}\right\rceil.
\]
\end{lemma}
\begin{proof}
Write the triple as $E,F,G$, with $X=E\cap G$, $Y=E\cap F$, $Z=F\cap G$, and relabel so $x=|X|\ge y=|Y|\ge z=|Z|$. Choose $P_G\subseteq G\setminus(E\cup F)$ of size $B-x-y$, and $B_0\subseteq F\setminus(E\cup G)$ of size $\min(B-x,k-y-z)$. Put $C=F\setminus(Y\cup B_0)$; then
\[
c=|C|=\max(k-B+x-y,z).
\]
Choose $P_E\subseteq E\setminus(F\cup G)$ of size $\min(B-x-c,k-x-y)$. All requested sizes are nonnegative and fit their hosts: $B-x-y\ge B-2m\ge0$, $c\le k-B+m$, $B-x-c\ge2B-k-2m\ge0$, and $B-x-y\le k-x-z$ because $y\ge z$ and $B\le k$.

Request edges avoiding
\[
X\cup Y\cup P_G,\quad X\cup B_0,\quad X\cup C\cup P_E,\quad (E\cup G)\setminus(X\cup P_E\cup P_G).
\]
The first three sizes are at most $B$. The fourth size is
\[
\max\{3(k-B)+x,\ 2(k-B)+y+z,\ (k-B)+2y\}\le B,
\]
by the definition of $B$.
A candidate pair with one point in $X$ has its other point in $Y,B_0$, or $C$ and is missed by one of the first three responses. Otherwise a point in $Y$ must pair with a point in $G\setminus X$, and the first or fourth response misses the pair according as the latter point is in $P_G$ or its complement. If the first point is in $Z$, use the third or fourth response and $P_E$. This exhausts the candidate pairs.
\end{proof}

\begin{lemma}[Splitting one pair cell]\label{lem:split}
For a good triple with pair sizes $x=|E\cap F|$, $y=|E\cap G|$, $z=|F\cap G|$, four legal responses yield a bad seven-tuple at any budget $T\le k$ satisfying
\[
T\ge\max\{(k+x+y+z)/2,\ k-x+|y-z|,\ k/3+x\}.
\]
The lemma is exactly integral.
\end{lemma}
\begin{proof}
Put $S=x+y+z$; the first hypothesis and $T\le k$ give $S\le T$. Request $H$ avoiding $X=E\cap F$, $Y=E\cap G$, $Z=F\cap G$, and
\[
p=\max(0,k-2(T-x)-y-z)
\]
points of the private part of $G$. This selection fits because $T\ge x$. Its union size is
\[
S+p=x+\max(y+z,k-2T+2x)\le T.
\]
All four triple cells vanish. Put $A=F\cap H$, $B=G\cap H$, $C=E\cap H$ with sizes $a,b,c$. Then
\[
a\le k-x-z,\quad b\le2(T-x),\quad c\le k-x-y,\quad a+b+c\le k.
\]

If $b\le T-x$, the three requests $X\cup B$, $Y\cup A$, $Z\cup C$ have sizes at most $T$: for the latter two use $k-x+|y-z|\le T$. They cover the three complementary candidate pairs.

If $b>T-x$, divide $B$ into two parts of size at most $T-x$; this is possible because $b\le2(T-x)$, also integrally. Two requests contain $X$ and one part each. The third contains $Y\cup Z\cup A\cup C$, whose size satisfies
\[
y+z+a+c\le k-b+y+z<k-T+x+y+z\le T.
\]
These requests again cover all three candidate pairs.
\end{proof}

\begin{lemma}[One surviving triple cell]\label{lem:fourcase}
Let $E,F,G$ be a good triple of $k$-sets, and put $X=E\cap F$, $Y=E\cap G$, $Z=F\cap G$, with sizes $x,y,z$ and $S=x+y+z$. Four further avoidance responses force a bad seven-tuple at any integer budget $T\le k$ satisfying
\[
\begin{gathered}
T\ge\max\Bigl\{x+y,\ k/2+x,\ k/2+y,\ k+x-y-z,\ k-x+y-z,\\
 k-S/3,\ (3k+S)/5,\ (k+x+y+2z)/3,\ (2k+3z)/4\Bigr\}.
\end{gathered}
\]
The construction is exactly integral. Any of the three pair cells can be designated $Z$.
\end{lemma}
\begin{proof}
Request $H$ avoiding $X,Y$ and a subset $Z_0\subseteq Z$ of size $\min(z,T-x-y)$. This is legal. Write
\[
Q=Z\cap H,\quad A=(F\cap H)\setminus Q,\quad B=(G\cap H)\setminus Q,\quad C=E\cap H,
\]
with sizes $q,a,b,c$. Also put $V=Z\setminus Q$, $v=z-q$, and $D=E\setminus(X\cup Y\cup C)$, $d=k-x-y-c$. These eight cells are disjoint, and
\[
q\le(S-T)_+,\quad a\le k-x-z,\quad b\le k-y-z,\quad c\le k-x-y,\quad q+a+b+c\le k.
\]
The only triple cell of $E,F,G,H$ is $Q$. A pair piercing these four edges therefore either belongs to $Q\times E$, or to one of $X\times B$, $Y\times A$, $V\times C$. We will make all three remaining avoidance requests contain $Q$ and their union contain $E$, while covering these last three products.

Two uniform bounds used below are
\[
q+x+b\le\max(k+x-y-z,k+2x-T)\le T,
\]
\[
q+y+a\le\max(k-x+y-z,k+2y-T)\le T.
\]
Also $z\le T$, since $4T\ge2k+3z$ and $z\le k$.

\paragraph{Case 0: $c\le T-z$.} Start the three requests with the disjoint-cell unions
\[
Q\cup X\cup B,\qquad Q\cup Y\cup A,\qquad Z\cup C.
\]
Each base has size at most $T$. Their total load, plus $d$, is
\[
k+2q+a+b+z\le\max(3k-S,3k+S-2T)\le3T.
\]
The first inequality follows from $a+b\le2k-x-y-2z$ and $q\le(S-T)_+$. Distribute $D$ among the remaining capacities. These capacities and $d$ are integers; hence the distribution can be integral. All candidate pairs are now covered by the bases or by $Q$ together with this distribution.

In the other cases $c>T-z$. Set
\[
a_0=k+x+z-2T,\qquad b_0=k+y+z-2T.
\]

\textbf{Case 1: $a<a_0$.} Use bases
\[
Q\cup X\cup B,\qquad Y\cup A\cup Z\cup C_2,\qquad Z\cup C_3,
\]
where $C=C_2\sqcup C_3$. This split fits because
\[
y+a+z\le k+x+y+2z-2T\le T,
\]
\[
a+c\le2k+z-y-2T\le2T-2z-y.
\]
Thus the two capacities $T-y-a-z$ and $T-z$ are nonnegative and sum to at least $c$. The total base load plus $d$ is
\[
k+q+a+b+2z\le2k+2z-c<2k+3z-T\le3T.
\]
Distribute $D$ among the remaining capacities. The three products are covered, and all bases contain $Q$.

\textbf{Case 2: $b<b_0$.} Interchange $X,B,x,b$ with $Y,A,y,a$ in Case 1. This covers the case even if both inequalities hold.

\textbf{Case 3: $a\ge a_0$ and $b\ge b_0$.} Put $u=c+z-T>0$. Since $z\le T$, we have $u\le c$, so choose $C_{12}\subseteq C$ of size $u$ and put $C_3=C\setminus C_{12}$. Choose an integer $t$ in
\[
\max(0,y+a+z+u-T)\ \le t\le\ \min(v,T-q-x-b-u).
\]
This interval is nonempty. The second upper entry is nonnegative because
\[
q+b+c\le k-a\le2T-x-z.
\]
The second lower entry is at most $v$ because
\[
q+a+c\le k-b\le2T-y-z.
\]
Finally the untruncated lower entry does not exceed the untruncated upper entry: this is equivalent to
\[
q+a+b+2c+x+y+3z\le4T,
\]
which follows from $q+a+b+c\le k$, $c\le k-x-y$ and $2k+3z\le4T$. All endpoints are integers.

Split $V=V_{13}\sqcup V_{23}$ with $|V_{13}|=t$, and use bases
\[
Q\cup X\cup B\cup C_{12}\cup V_{13},
\]
\[
Q\cup Y\cup A\cup C_{12}\cup V_{23},
\]
\[
Z\cup C_3.
\]
The interval inequalities make the first two sizes at most $T$; the third has size exactly $T$. Their total load plus $d$ is
\[
k+q+a+b+2z+u\le2k+3z-T\le3T.
\]
Again distribute $D$ into the residual capacities. The first two requests cover $X\times B$, $Y\times A$ and $V\times C_{12}$; the third covers $V\times C_3$. All contain $Q$ and their union contains $E$. Thus each candidate piercing pair of the first four edges misses at least one of the three new responses. This is a bad seven-tuple.
\end{proof}

\begin{lemma}[An asymmetric split]\label{lem:asymone}
With the good-triple notation of Lemma~\ref{lem:fourcase}, four further avoidance responses force a bad seven-tuple whenever $T\le k$ and
\[
T\ge\max\left\{S,\ \frac k2+y,\ \frac{k+2x-y+z}{2},\ \frac{k+2x+y+3z}{3}\right\}.
\]
The construction is exactly integral. All six permutations of $x,y,z$ are available.
\end{lemma}
\begin{proof}
Request $H$ avoiding $X,Y,Z$ and $T-S$ points in the private part of $F$. The private selection fits because $T-S\ge0$ and $T-S\le k-x-z$, the latter following from $T\le k+y$. Put $A=F\cap H$, $B=G\cap H$, $C=E\cap H$, with sizes $a,b,c$. All four triple cells vanish, so the candidate piercing pairs are $X\times B$, $Y\times A$, and $Z\times C$. Also
\[
a\le k-T+y,\qquad b\le k-y-z,\qquad a+b+c\le k.
\]
If $a\le T-y-z$, split $B=B_1\sqcup B_2$ into two parts of size at most $T-x-z$. This is possible, integrally, since
\[
2(T-x-z)\ge k-y-z\ge b.
\]
Start the final three requests with
\[
X\cup Z\cup B_1,\qquad X\cup Z\cup B_2,\qquad Y\cup Z\cup A.
\]
The first two fit by construction; the third fits by the case assumption. The total base load plus $c$ is
\[
2x+y+3z+a+b+c\le k+2x+y+3z\le3T.
\]
Distribute $C$ integrally among the residual capacities. Every final request contains $Z$, so this distribution covers $Z\times C$; the first two requests cover $X\times B$ and the third covers $Y\times A$.

If instead $a>T-y-z$, then
\[
b+c\le k-a<k-T+y+z\le2(T-x-z).
\]
Split $B\cup C$ into two parts of size at most $T-x-z$, and put $X\cup Z$ together with one part in each of the first two requests. The third request is $Y\cup A$, of size at most $k-T+2y\le T$. These three requests cover all three candidate-pair products again. Thus no pair pierces the seven edges in either case.
\end{proof}

\begin{lemma}[A second asymmetric split]\label{lem:asymtwo}
With the same good-triple notation, four further avoidance responses force a bad seven-tuple if $T\le k$ and
\[
T\ge\max\left\{S,\ \frac k3+x,\ k-x+z,\ \frac{k+2x+3y+z}{3}\right\}.
\]
The construction is exactly integral, with all six permutations available.
\end{lemma}
\begin{proof}
Request $H$ avoiding $X,Y,Z$ and
\[
p=\max(0,k-2(T-x)-y-z)
\]
points of the private part of $G$. This fits its host since $T\ge x$, and its total cost is $\max(S,k+3x-2T)\le T$. As in Lemma~\ref{lem:split}, write $A=F\cap H$, $B=G\cap H$, $C=E\cap H$, with sizes $a,b,c$. Then
\[
b\le2(T-x),\qquad c\le k-x-y,\qquad a+b+c\le k.
\]
All four triple cells vanish and the candidate pairs are $X\times B$, $Y\times A$, $Z\times C$.

If $b>k+y+z-T$, split $B$ into two parts of size at most $T-x$, and use $X$ with each part as two requests. The third is $Y\cup Z\cup A\cup C$, of size at most $k-b+y+z<T$.

Otherwise $b\le k+y+z-T\le2(T-x-y)$, where the last inequality is exactly $3T\ge k+2x+3y+z$. Split $B=B_1\sqcup B_2$ into two parts of size at most $T-x-y$. Start the three requests with
\[
X\cup Y\cup B_1,\qquad X\cup Y\cup B_2,\qquad Y\cup Z\cup C.
\]
The first two fit by construction; the third has size at most $k-x+z\le T$. The total base load plus $a$ is at most
\[
2x+3y+z+a+b+c\le k+2x+3y+z\le3T.
\]
Distribute $A$ into the residual integer capacities. Since every request contains $Y$, this covers $Y\times A$. The first two cover $X\times B$ and the third covers $Z\times C$. Both cases therefore yield a bad seven-tuple.
\end{proof}

\begin{lemma}[Two traces forced through a gap]\label{lem:gapempty}
Suppose no pair of family edges has intersection in $[\ell k,hk]$. For a good triple with pair sizes $x,y,z$, define $S=x+y+z$. If each of
\[
S,\quad (1-h)k+y,\quad (1-h)k+z,\quad (2-2h)k-x,
\]
\[
k/3+x,\quad ((2-h)k+2x-y)/3,\quad ((2-h)k+2x-z)/3,
\]
\[
((3-2h)k+x-y-z)/3,\quad 2\ell k+y+z
\]
is at most $\beta k$, and $T=\lceil\beta k\rceil+1\le k$, four further requests of size at most $T$ produce a bad seven-tuple. Here $x$ is the pair cell shared by the two edges whose new traces will be forced below $\ell k$.
\end{lemma}
\begin{proof}
Write the edges as $E,F,G$, with $X=E\cap F$, $Y=E\cap G$, $Z=F\cap G$. Let $h_0=\lfloor hk\rfloor$ and choose private subsets of $E,F$ of sizes
\[
p=(k-x-y-h_0)_+,\qquad t=(k-x-z-h_0)_+.
\]
Avoid them together with $X,Y,Z$. The cost before padding is
\[
C_0=x+\max(y,k-x-h_0)+\max(z,k-x-h_0).
\]
Expanding these two maxima gives the first four displayed forms, with rounding error strictly less than two. Thus $C_0<\beta k+2$, and its integrality gives $C_0\le\lceil\beta k\rceil+1=T$. Pad the request with $T-C_0$ private points of $G$. This fits: $C_0\ge x+y+z$ and $T\le k$ imply $T-C_0\le k-y-z$.

The response $H$ has its $E$- and $F$-traces at most $h_0$, so the gap makes both strictly smaller than $\ell k$. Its $G$-trace has size $b\le k-T+x+p+t$. All triple intersections among $E,F,G,H$ are empty. Split the $G$-trace into two nearly equal integer parts and use $X$ together with each part in two requests. Their sizes are at most
\[
\frac{k+3x+p+t-T}{2}+\frac12.
\]
Expanding the two positive parts in $p,t$, the numerator $k+3x+p+t$ is the maximum of
\[
k+3x,\quad (2-h)k+2x-y,\quad (2-h)k+2x-z,\quad (3-2h)k+x-y-z,
\]
with error strictly less than two. The next four hypotheses therefore make each split request strictly less than $(3\beta k+2-T)/2+1/2\le T$, since $T\ge\beta k+1$.

Use $Y,Z$ and the two small traces of $H$ in the third request. Its size is strictly smaller than $y+z+2\ell k\le\beta k$. These three requests eliminate the three complementary candidate-pair products.
\end{proof}

\begin{lemma}[A surviving triple cell and a gap]\label{lem:gapsurvive}
Suppose no pair of family edges has intersection in $[\ell k,hk]$. For a good triple with the usual pair sizes $x,y,z$, put $S=x+y+z$. Suppose $S\ge\beta k$, and each of
\[
x+y,\quad (1-h)k+y,\quad (1-h)k+z,
\]
\[
k/4+x+(y+z)/4,\quad 2\ell k+y+z,\quad k/2+x/2+(z+\ell k)/4
\]
is at most $\beta k$. If $T=\lceil\beta k\rceil+1\le k$, four further requests of size at most $T$ force a bad seven-tuple. The assertion includes the rounding strip $\beta k\le S<T$.
\end{lemma}
\begin{proof}
Request $H$ avoiding $X,Y$ and a subset $Z_0\subseteq Z$ of size $\min(z,T-x-y)$. This fits its budget since $x+y\le\beta k\le T$. Put $Q=Z\cap H$, $A=(F\cap H)\setminus Q$, $B=(G\cap H)\setminus Q$, $C=E\cap H$, with sizes $q,a,b,c$. Then
\[
q\le(S-T)_+\le S-\beta k,\qquad b\le k-y-z.
\]
Moreover
\[
c\le k-x-y=k-S+z\le(1-\beta)k+z\le hk,
\]
\[
a+q\le k-x-z+(S-T)_+\le(1-\beta)k+y\le hk.
\]
For the second inequality, use $S\ge\beta k$ if $S\le T$, and $T\ge\beta k$ otherwise. These are actual pair-intersection sizes, so the gap forces $c<\ell k$ and $a+q<\ell k$.

Split $B=B_1\sqcup B_2$ into nearly equal integer parts. Start two requests with $X\cup Q\cup B_1$ and $X\cup Q\cup B_2$. Their sizes are at most
\[
x+q+\lceil b/2\rceil\le2x+(y+z)/2+k/2-\beta k+1/2\le\beta k+1/2\le T.
\]
Start the third request with $Y\cup Z\cup A\cup C$, of size strictly below $y+z+2\ell k-q\le\beta k$.

All three bases contain $Q$. To make their union contain all of $E$, distribute $D=E\setminus(X\cup Y\cup C)$ among the residual capacities. The total base load plus $|D|$ is
\[
k+x+z+2q+a+b<k+x+z+q+\ell k+b
\]
\[
\le2k+x-y+\ell k+q\le2k+2x+z+\ell k-\beta k\le3\beta k\le3T.
\]
The distribution therefore exists with integer sizes. Candidate piercing pairs of $E,F,G,H$ are in $Q\times E$, $X\times B$, $Y\times A$, or $(Z\setminus Q)\times C$. The first kind is eliminated because every final request contains $Q$ and their union contains $E$; the other three kinds are eliminated by the indicated bases.
\end{proof}

\begin{lemma}[Two large opposite pair cells]\label{lem:twolarge}
Suppose all pair intersections in the family are at most $m$ or greater than $k/2$, where $m\le k/4$. Let $E,F,G,H$ be four $k$-edges with all triple intersections empty. Put $X=E\cap F$, $B=G\cap H$, of sizes $x,b>k/2$, and suppose the other four pair cells each have size at most $m$. Three further legal responses give a bad seven-tuple if $T\le k$ and
\[
T\ge\lceil k/2\rceil+2m,\qquad T\ge x+m,\qquad 3T\ge2x+b+\lceil k/2\rceil+2m.
\]
\end{lemma}
\begin{proof}
Put $Y=E\cap G$, $Z=F\cap G$, $A=F\cap H$, $C=E\cap H$, and let $U=Y\cup Z\cup A\cup C$. These six pair cells are disjoint. Write $s=|Y|+|Z|$, $t=|A|+|C|$, so $s,t\le2m$, and define
\[
p=\lceil k/2\rceil-\min(s,t),\qquad q=T-\lceil k/2\rceil-\max(s,t).
\]
Both are nonnegative. Moreover $p\le\lceil k/2\rceil\le b$ and $q\le T-\lceil k/2\rceil\le\lfloor k/2\rfloor<x$. Choose $B_0\subseteq B$, $X_0\subseteq X$ with sizes $p,q$, and request $I$ avoiding $U\cup B_0\cup X_0$, whose size is exactly $T$.

The $G$-trace of $I$ has size at most
\[
k-s-p=\lfloor k/2\rfloor+\min(s,t)-s\le\lfloor k/2\rfloor,
\]
and the analogous statement holds for $H$. The global gap therefore forces both traces to have size at most $m$. In particular $d=|B\cap I|\le m$. Also $x_I=|X\cap I|\le x-q$.

Choose $B_1\subseteq B$ containing $B\cap I$ and having size $\min(b,T-x)$. This is possible since $T-x\ge m\ge d$. Use the last two requests
\[
X\cup B_1,\qquad (X\cap I)\cup(B\setminus B_1).
\]
The first has size at most $T$. For the second, note that
\[
x+x_I+b\le2x+b-q\le2x+b-T+\lceil k/2\rceil+2m\le2T.
\]
It follows that its size $x_I+\max(0,b-T+x)$ is at most $T$ as well.

A pair piercing $E,F,G,H$ belongs to $X\times B$, $Y\times A$ or $Z\times C$. The latter two products lie entirely in $U$ and hence miss $I$. For a remaining pair in $X\times B$, if its $B$-point belongs to $B_1$, the penultimate response misses both points. Otherwise its $B$-point lies outside $I$, since $B\cap I\subseteq B_1$; its $X$-point must then belong to $I$, and the last response misses both. This proves the lemma.
\end{proof}

\begin{lemma}[Two surviving triple cells]\label{lem:gaptwo}
Suppose no pair intersection lies in $[\ell k,hk]$. For a good triple $E,F,G$ let $X=E\cap F$, $Y=E\cap G$, $Z=F\cap G$, with sizes $x,y,z$, and put $S=x+y+z$. Suppose
\[
x+y\ge(1-h)k,\qquad y+z\ge(1-h)k,
\]
and each of
\[
(2-2h)k-y,\quad x+\ell k,\quad z+\ell k,\quad k/2+y,\quad 3k/4,\quad (3k+S)/5
\]
is at most $\beta k$. Then four further responses of avoidance budget $T=\lceil\beta k\rceil+1\le k$ give a bad seven-tuple. All six orientations are allowed.
\end{lemma}
\begin{proof}
Let $h_0=\lfloor hk\rfloor$ and $g=(k-y-h_0)_+$. The domain inequalities imply $g\le x,z$. Choose $g$ points from each of $X,Z$, and avoid them together with all of $Y$. The initial cost is
\[
y+2g=\max(y,2k-y-2h_0)<\beta k+2.
\]
Since the cost is an integer, it is at most $\lceil\beta k\rceil+1=T$. Complete the portion of the request inside $F$ to exactly $T-y$ points, first using any remaining points of $X\cup Z$ and then private points of $F$. This is possible because $X,Z$ are disjoint subsets of $F$, $Y$ is disjoint from $F$, and $T-y\le k$. Let $H$ be an avoiding response.

Write $P=E\cap F\cap H$ and $Q=F\cap G\cap H$, of sizes $p,q$. No point of $Y$ is in $H$, so these are the only possible triple cells among the four edges. Set
\[
A=(F\cap H)\setminus(P\cup Q),\quad B=(G\cap H)\setminus Q,\quad C=(E\cap H)\setminus P,
\]
with sizes $a,b,c$. The priority given to $X\cup Z$ in the first request yields
\[
p+q\le(S-T)_+,
\]
and its total size inside $F$ gives
\[
d:=|F\cap H|=p+q+a\le k-T+y.
\]
The two $g$-point cuts also give $|E\cap H|,|G\cap H|\le h_0$. The gap therefore makes both intersections smaller than $\ell k$.

Start the final three requests with
\[
X\cup(G\cap H),\qquad Z\cup(E\cap H),\qquad Y\cup(F\cap H).
\]
Their sizes are at most $x+\ell k$, $z+\ell k$, and $k-T+2y$, respectively, so all fit in $T$. Each contains both $P$ and $Q$.

Let $D=E\setminus(F\cup G\cup H)$ and $W=G\setminus(E\cup F\cup H)$. These are disjoint, with sizes $k-x-y-c$ and $k-y-z-b$. The sum of the three base sizes plus $|D|+|W|$ is
\[
2k-y+2(p+q)+a=2k-y+(p+q)+d\le3k-T+(S-T)_+.
\]
This is at most $3T$: if $S\le T$, use $T\ge3k/4$; otherwise use $5T\ge3k+S$. Consequently $D\cup W$ can be partitioned among the residual integer capacities of the three requests.

Their union now contains $E\cup G$. A candidate piercing pair involving $P$ has its other point in $G$, while a pair involving $Q$ has its other point in $E$; all such pairs are eliminated since every request contains $P,Q$ and their union contains $E,G$. Every other candidate pair lies in $(X\setminus P)\times B$, $(Z\setminus Q)\times C$, or $Y\times A$, contained in the first, second, or third request respectively. Their three avoiding responses give the contradiction.
\end{proof}

\section{Finishing when the largest small intersection is small}\label{app:smallmax}
\begin{lemma}\label{lem:smallmax}
Let $T$ be an integer with $6k/7\le T\le k$, satisfying
\begin{equation}\label{eq:fourinteger}
 \begin{gathered}
 6T>5k+1,\qquad 3T\ge2k+\ceil{k/2},\\
 2T\ge3k/2+1/2,\qquad 12T\ge10k+4.
 \end{gathered}
\end{equation}
If every pair intersection in a $k$-uniform hypergraph with property
$(7,2)$ is at most $2k/7$ or greater than $k/2$, then $\tau(\HH)\le T$.
\end{lemma}
\begin{proof}
Suppose $\tau(\HH)>T$. Choosing an edge avoiding a $T$-subset of
another gives an intersection at most $k-T$. Thus the largest pair
intersection $m\le k/2$ exists, and $m\le2k/7$. Choose a pair attaining
$m$ and use Lemma~\ref{lem:balanced}. Both new intersections are at most
$k-\lfloor(T+m)/2\rfloor$.

If both are at most $k/2$, then all three are at most $m$ by maximality.
Lemma~\ref{lem:smalltriple} applies at a budget at most $T$, since
\[
 (2k+2m)/3\le6k/7,\qquad (3k+m)/4\le23k/28<6k/7.
\]

Otherwise, relabel the good triple so that
\[
 x=|E\cap F|>k/2,\qquad |E\cap G|=m,\qquad |F\cap G|\le m.
\]
At most one of the new intersections is greater than $k/2$, since they
are disjoint subsets of the response edge. The balanced bound gives
\begin{equation}\label{eq:smallmaxbounds}
 x\le k-(T+m)/2+1/2,\qquad m\le k-T.
\end{equation}
For the second assertion, $T+m\ge k+1$ would imply
$\lfloor(T+m)/2\rfloor\ge\lceil k/2\rceil$, contrary to $x>k/2$.
Now $T\ge6k/7$ gives $m\le k/7<k/4$, and
\[
 x+2m\le(5k-4T+1)/2<T,
\]
since $6T>5k+1$. Request $H$ avoiding the three pair cells and enough
private points of $G$ to make the request size $T$. The padding fits:
the initial request size is $x+|E\cap G|+|F\cap G|$, and $T\le k$.
All triple intersections among $E,F,G,H$ are now empty. The $E$- and
$F$-traces of $H$ are less than $k/2$, because $x>k/2$, and are
therefore at most $m$. Write
\[
 b=|G\cap H|\le k-T+x.
\]
If $b\le k/2$, apply Lemma~\ref{lem:smalltriple} to $E,G,H$.

If $b>k/2$, the four edges satisfy the hypotheses of
Lemma~\ref{lem:twolarge}. Its three budget requirements follow from
\begin{align*}
 \ceil{k/2}+2m
   &\le\ceil{k/2}+2k-2T\le T,\\
 x+m&\le3k/2-T+1/2\le T,\\
 2x+b+\ceil{k/2}+2m
   &\le(9k+m-5T+4)/2\le3T.
\end{align*}
For the last inequality use $m\le k-T$ and $12T\ge10k+4$.
The first two use the second and third inequalities of
\eqref{eq:fourinteger}. The lemma produces a bad subfamily of at most seven
edges, completing the contradiction.
\end{proof}


\section*{Research assistance}
Claude (Anthropic) and Codex (OpenAI) were used for literature searches,
mathematical exploration, computational experiments, and manuscript
preparation. The proofs presented here are self-contained hand arguments;
no numerical output or software execution is required to verify them.

\begin{thebibliography}{9}
\bibitem{EFKT}
P.~Erd\H{o}s, D.~G.~Fon-Der-Flaass, A.~V.~Kostochka and Zs.~Tuza,
\emph{Small transversals in uniform hypergraphs},
Siberian Advances in Mathematics \textbf{2} (1992), 82--88.

\bibitem{FKW}
D.~G.~Fon-Der-Flaass, A.~V.~Kostochka and D.~R.~Woodall,
\emph{Transversals in uniform hypergraphs with property $(7,2)$},
Discrete Mathematics \textbf{207} (1999), 277--284.
\href{https://doi.org/10.1016/S0012-365X(99)00114-4}
{doi:10.1016/S0012-365X(99)00114-4}.
\bibitem{Kostochka}
A.~V.~Kostochka,
\emph{Transversals in uniform hypergraphs with property $(p,2)$},
Combinatorica \textbf{22} (2002), 275--285.
\href{https://doi.org/10.1007/s004930200013}{doi:10.1007/s004930200013}.

\bibitem{BKS}
M.~Buci\'c, D.~Kor\'andi and B.~Sudakov,
\emph{Covering graphs by monochromatic trees and Helly-type results
for hypergraphs}, Combinatorica \textbf{41} (2021), 319--352.
\href{https://doi.org/10.1007/s00493-020-4292-9}
{doi:10.1007/s00493-020-4292-9}.
\end{thebibliography}
\end{document}
